\chapter{Introduction}

\section{Context}
The General Data Protection Regulation (GDPR) requires organizations to maintain a Record of Processing Activities (ROPA), a formal document that outlines how personal data is collected, processed, stored, and shared. Creating and maintaining these ROPA documents is often a resource-intensive process, especially for small and medium-sized organizations that lack specialized legal or data-protection expertise.
In response to this challenge, ROPAgen --- the tool evaluated in this thesis --- is a web-based application that uses Large Language Models (LLMs) to assist users in drafting GDPR-compliant ROPA documentation. ROPAgen was developed by the author at the Institute of Databases and Information Systems (DBIS), University of Ulm, prior to this evaluation; it provides three interaction modes (Form, Ask, and Chat) that together offer increasing levels of LLM support, described in detail in \S\ref{sec:ropagen}. The system aims to reduce the effort needed to create legally compliant documentation while maintaining transparency and user control.
Although the application provides a functional framework for assisted ROPA generation, its effectiveness, usability, and impact on documentation quality have not yet been systematically evaluated. To ensure that such a system can be confidently integrated into practical compliance workflows, both qualitative and quantitative assessments are needed. These evaluations must determine whether users can efficiently and correctly create ROPA documents, how the varying levels of LLM support influence performance and trust, and whether the system produces outputs aligned with GDPR requirements.

 
\section{Research Problem and Research Questions}
The goal of this thesis is to conduct a comprehensive evaluation of the existing LLM-supported ROPA generation application. The evaluation asks how the three interaction modes shape the experience of novice users, how they shape the documents those users produce, and whether the two are aligned. Each overarching research question is broken into three more specific sub-questions tied to the instruments used to answer them.

\paragraph{RQ1 --- User experience.}
\emph{How do varying degrees of LLM support influence novice users during ROPA document creation?}
\begin{itemize}
    \item[\textbf{RQ1.1}] How do novice users perceive the \emph{usability} of the three interaction modes (Form, Ask, Chat)?
    \item[\textbf{RQ1.2}] How does \emph{cognitive load} differ across the three modes for users with no prior GDPR knowledge?
    \item[\textbf{RQ1.3}] How do novice users subjectively evaluate each mode --- in their perception of the AI's support, in their confidence about the documents they produce, and in the mode they ultimately prefer?
\end{itemize}

\paragraph{RQ2 --- Document quality.}
\emph{How do varying degrees of LLM support influence the quality of ROPA documents produced by novice users?}
\begin{itemize}
    \item[\textbf{RQ2.1}] How does document quality differ across the three modes when measured by reference-based NLP similarity to an expert-authored ROPA (BLEU, ROUGE, METEOR, BERTScore, SBERT)?
    \item[\textbf{RQ2.2}] As a supplementary, reference-free signal, what additional GDPR-specific quality properties (completeness, scenario faithfulness, legal correctness, hallucination) does an LLM-as-a-judge surface across the three modes?
    \item[\textbf{RQ2.3}] Is there a discrepancy between user confidence (RQ1.3) and the document quality actually achieved (RQ2.1, RQ2.2)?
\end{itemize}
